% =========================================================
%  SIREN -- Related Work
%  Target: Machine Learning (Springer, Q1)
% =========================================================

\section{Related Work}
\label{sec:related}

\subsection{Machine Learning for Intrusion Detection}

Early approaches to network intrusion detection relied heavily on traditional machine learning techniques. \citet{dang2019studying} provided a comprehensive survey of these methods, demonstrating that tree-based ensemble learning combined with careful feature engineering could achieve state-of-the-art performance on tabular flow data. However, they also highlighted a critical limitation: modern machine learning algorithms are extremely data-hungry and require exhaustive labelling efforts, motivating the search for approaches that can operate with minimal or no labelled attack data. This limitation underscores the necessity for the unsupervised and structure-aware models developed in recent years. Furthermore, to address severe class imbalances and improve robustness in tabular models, \citet{dang2026utilizing} recently proposed integrating uncertainty measures and isolation forests, significantly reducing false negative rates for minority attack classes. While effective for flat vector spaces, these uncertainty-aware paradigms have yet to be natively integrated into dynamic graph structures.

\subsection{Graph Neural Networks for Network Intrusion Detection}

Graph neural networks (GNNs) have become the dominant paradigm for network
intrusion detection owing to their capacity to exploit the structural
relationships among hosts and traffic flows that elude flat feature-vector
classifiers~\citep{Zhong2024survey}.
\citet{Lo2022egraphsage} proposed E-GraphSAGE, converting per-flow records
into a network graph and applying GraphSAGE-based edge classification;
despite its widespread use as a baseline, E-GraphSAGE operates on static
snapshot graphs and emits uncalibrated point predictions.
More recent work has sought to enrich both graph construction and
representation learning.
\citet{Ngo2025tksgf} introduced TKSGF, replacing physical-link graphs with
Top-$K$ attribute-similarity edges paired with a GraphSAGE backbone, and
demonstrated that feature-driven topology substantially improves detection
accuracy on IoT benchmarks.
\citet{Lin2025egracl} proposed E-GRACL, augmenting GraphSAGE with residual
connections, a global attention mechanism, local gating, and graph
contrastive learning to capture both edge features and topological information
simultaneously.
At the multi-scale level, \citet{Liu2025mgfgnn} built MGF-GNN, fusing
host-level, port-level, and protocol-level graphs inside a hierarchical
message-passing framework, while \citet{Zhang2026aedgnn} addressed adversarial
evasion through AEDGNN, a semi-supervised GNN formulation that explicitly
models inter-flow relationships to resist labelling noise.
Parallel work by \citet{Ahanger2025gat} applied graph attention networks to
IoT intrusion detection, and \citet{VillegasCh2025dynamic} developed a
dynamic graph modelling approach for IoT traffic.
All of these methods share a common limitation with respect to \textsc{Siren}:
they operate on discrete-time graph snapshots, produce point predictions, and
provide no probabilistic guarantee over the \emph{stationary behaviour} of the
monitored system under normal conditions.

\subsection{Score-Based Generative Models}

The score function $\nabla_{\mathbf{x}} \log p(\mathbf{x})$ was introduced
as a tractable surrogate for the intractable log-likelihood by
\citet{Hyvarinen2005score}, who proved that minimising the expected squared
difference between the model score and the data score---score
matching---equals maximum likelihood under mild regularity conditions.
\citet{Vincent2011dsm} later showed that denoising autoencoders implicitly
perform score matching, motivating the denoising score matching (DSM)
objective that sidesteps the intractable Hessian.
Building on this, \citet{Song2019ncsn} proposed NCSN, training a
noise-conditional score network across a ladder of noise levels and using
Langevin dynamics at inference; this multi-scale strategy dramatically
improves score estimation quality in high-dimensional settings.
\citet{Song2021sde} then unified diffusion-based and score-based
generative models under a continuous-time SDE framework, in which the
forward process is a prescribed It\^o SDE that injects noise and the learned
reverse drift is parameterised by the score network.
\textsc{Siren} borrows the multi-scale DSM objective and the score-SDE
connection from these works, but redirects them toward a fundamentally
different objective: rather than generation, we exploit the learned score of
the \emph{stationary distribution} of a graph-coupled SDE as a structured
anomaly signal.

\subsection{Continuous-Time Models on Graphs}

Static GNNs~\citep{Kipf2017gcn,Hamilton2017graphsage} capture structural
relationships at a single snapshot but are inherently unable to model the
continuous temporal evolution of network traffic.
Temporal graph networks (TGN)~\citep{Rossi2020tgn} extended GNNs to dynamic
graphs through memory modules and temporal attention, enabling event-level
updates; however, TGN operates in discrete event time and does not posit a
generative stochastic process over the joint graph state.
Graph neural diffusion (GRAND)~\citep{Chamberlain2021grand}
re-interpreted message passing as the discretisation of a \emph{deterministic}
heat-diffusion partial differential equation on the graph, providing improved
over-smoothing control and depth generalisation; yet the deterministic ODE
formulation cannot model the inherent stochasticity of real network traffic.
Closest in spirit to our approach, \citet{Kidger2021nsde} demonstrated that
neural stochastic differential equations can be trained end-to-end via
adjoint-based backpropagation through the It\^o integral, establishing the
computational feasibility of learned stochastic dynamics on sequences.
\textsc{Siren} extends this programme to \emph{coupled} systems of It\^o SDEs
on graphs, where the graph Laplacian enters the drift, the diffusion is
learned per-node, and the stationary distribution of the resulting process
provides the detection signal---an objective and architecture that is absent
from all prior graph--SDE work.

\subsection{Stein Operators and Graph Anomaly Detection}

The Stein operator $\mathcal{A}_{p}\,g(\mathbf{x})
= \nabla_{\mathbf{x}}\!\log p(\mathbf{x})^\top g(\mathbf{x})
  + \nabla_{\mathbf{x}}\!\cdot\! g(\mathbf{x})$
satisfies $\mathbb{E}_{p}[\mathcal{A}_{p}\,g] = 0$ for all smooth test
functions $g$ under mild regularity conditions.
\citet{Liu2016ksd} exploited this Stein identity to construct the kernelised
Stein discrepancy (KSD), a tractable goodness-of-fit measure that does not
require normalisation of the reference density---precisely the property that
makes it applicable when $p = \pi^{*}$ is known only through its score.
In the graph anomaly detection literature, \citet{Cai2021strgnn} proposed
StrGNN, a structural-temporal GNN that detects anomalous edges by mining
unusual temporal subgraph structures; applied to enterprise security, it
achieved zero false negatives on benchmark datasets.
Most recently, \citet{Bai2025crcsgad} introduced CRC-SGAD, applying
conformal risk control to supervised graph anomaly detection with formal
bounds on both the false-negative rate and the false-positive rate via a
spectral GNN calibrator.
CRC-SGAD is the closest existing work to \textsc{Siren} in combining
GNN-based detection with statistical guarantees for a security application;
however, it relies on conformal prediction---which requires labelled
calibration data---and does not posit or exploit a stochastic process model
for normal traffic.
\textsc{Siren} addresses this gap by grounding the anomaly signal in the
Stein identity of the \emph{learned} stationary distribution of a
graph-coupled It\^o SDE, requiring no attack labels at training or
calibration time.

\subsection{Distributionally Robust Learning on Graphs}

Distributionally robust optimisation (DRO) hardens models against
distributional shift by minimising the worst-case expected loss within a
Wasserstein ball around the empirical distribution~\citep{Sinha2018wrm}.
\citet{Sadeghi2021drognn} applied this principle to graph-structured data,
showing that a Wasserstein-ball DRO objective for GNN semi-supervised
classification admits a strong-dual reformulation as regularised empirical
risk minimisation.
\citet{Wang2024drgnn} extended graph DRO to collaborative filtering,
reinterpreting the GNN as a graph-smoothing regulariser to make the DRO
objective tractable at scale.
These methods improve worst-case robustness \emph{at training time} for
supervised tasks but leave inference uncalibrated and provide no anomaly
signal grounded in a generative process model.
\textsc{Siren} takes a complementary approach: rather than robust supervised
learning, it learns a continuous-time generative model of normal behaviour
and uses deviations from that model's stationary-distribution score as the
detection criterion---an unsupervised, process-model-based alternative that
requires no attack labels.